Induction on Lemma \claimref{sheaf.tame-wall-induction} via Remark \ref{rem:bundlesetup}, choosing at each stage a good coupling (possible: each $\Sig^{(k)}$ is a closed subset of a compact collar of its predecessor, hence compact, so its $S$ is finite, and only an upper bound on $|\mu\nu|$ is demanded). Each stage also needs its radical line trivial, so that collar data exist for the next coupling (Definition \claimref{obj.sheaf.collar-data-coupling}); this is automatic, because the shadow wall's radical line carries an explicit global spanning section: writing $\alpha_0 := G^{-1}e$ for the metric dual of the transversal functional, so that $s = \alpha_0(u)$, the section
\[
x \;:=\; \alpha_0 \,+\, \mu\beta s\, \nO \,+\, \nu\beta s\, \nInf
\]
satisfies $h(x, \alpha) = e(\alpha)\,(1 + 2\mu\nu\beta^2 s) = e(\alpha)\,F = 0$ on $\{F = 0\}$ for every $\alpha \in T^*M$, while $h(x, \nO) = \nu\beta s - \nu\beta s = 0$ and $h(x, \nInf) = \mu\beta s - \mu\beta s = 0$: so $x$ spans $\ker h$ along $\Sig'$, globally and without zeros. A global transversal for the next stage is then any $u'$ with $x(u') = 1$, available from any bundle metric. The ranks and signatures are Theorem \claimref{sheaf.one-sided-completion} at each stage. Figure \ref{fig:wallsequence} displays the sequence.
